

My name is Jordan Ling Shen Hang and I saw the listing for this position on Jobstreet, and I was instantly piqued. This is exactly the kind of position I've been seeking, an opportunity to work on a project from an analyst's perspective.

As a computer science student who has helped direct events in clubs at my university, I have found much enjoyment in planning something, and seeing all the pieces come together. My approach to planning events and other projects involves an iterative process of gathering the requirements by liaising with the relevant parties and building a plan, and then communicating the plan to gain feedback, and repeat. This helps me ensure that whatever is done best reflects the best interest of all the stakeholders involved. 

After the planning phase, I hold meetings to brief my team on the plan, and discuss the delegation and timeliness of tasks. Sometimes even when the plan seems solid, there is always something that is overlooked, and pointed out during this discussion. Which is why I find communicating with a team so important as they will see things that I may miss. 

% As mentioned on my resume I have experience in working in large events, namely CodeNection in 2025 (iFAST was a sponsor!) where I led a team of 5: regularly communicating with the directors to ensure alignment as well as mentors, judges, and participants.

I acknowledge that my approach to event planning is not perfectly representative for projects, and furthermore I would be an intern. But I believe that certain soft skills such as communication and effective planning are transferrable. Furthermore, with my background in computer science I also understand software engineering concepts, as well as data science concepts.

I am sure that if I am selected, this will be a very different experience than my experiences in university. But I am eager to learn, and I believe an internship at iFAST would be a great place for me to take my first steps into the tech industry.

Thank you for your time and consideration.


% Dear Maureen and hiring team,

% I was so excited to see your post on LinkedIn because it’s exactly the type of job I’m looking for: an opportunity to bring my experience with video production and enthusiasm for storytelling to an organization that sets the standard for high-quality management content.

% In addition to five years of experience in broadcast journalism, research, and video production, I would bring an organized and systems-level perspective to this role. I view video production as a puzzle, and like to think about which parts need to come together in order to make a great final product. My approach is to have in-depth conversations with my team members, and the various stakeholders, before each project. This helps me nail down the logistics — from location to talent.

% From there, the fun begins: fleshing out the concept and identifying what visuals will best represent it. Ideation and storyboarding are essential in this step. I know I’m not right all the time, so I enjoy working with a diverse team that can bring in new perspectives, brainstorm, and pitch ideas that will make the final product stronger. Whenever possible, I also try to seek out other sources for inspiration, like magazines, which allow me to observe different ways of expression and storytelling. This approach has served me well. It’s what has allowed me to enter the film industry and grow as a creator.

% On my website, you can see examples of how I use the above process to create fun, engaging content.

% Given this experience and my enthusiasm for the work you do, I believe I’d make a great addition to your team. I recently had a chance to try out your Patient Zero product at my current organization. The simulation is both challenging and engaging. I was impressed by your ability to apply  different storytelling methods to an online training course (which, let’s admit, can often be a little dry). Your work exemplifies exactly what I believe: There’s an opportunity to tell a compelling story in everything — all you have to do is deliver it right.

% I’d love to come in and speak with you more about what I’d be able to offer in this role. Harvard Business Publishing is my top choice and I believe I’d make valuable contributions to your team.

% Thank you for your time and consideration!